\section{Study Scripture and liturgical locators}
The Scripture printed here is the public-domain Douay--Rheims Bible, Challoner revision, in the retained Project Gutenberg 1581 witness. It is \textbf{study English, not the approved English proclaimed at this celebration}. The U.S. Lectionary for Mass, no. 139, controls the proclaimed readings. Original spelling is retained. A biblical basis printed for an adaptation does not reconstruct the Missal antiphon or the adapted acclamation. The orations are identified without substitute translations.

\subsection{Entrance}
\label{proper-entrance}
\properlocus{Cf. Esther 4:17; identified Douay--Rheims basis 13:9--11.}
The biblical prayer below is the basis of the adaptation, not its exact appointed wording. Missal identifier: \latin{In voluntate tua, Domine}.
\begin{englishwitness}{Douay--Rheims Challoner}
\textsuperscript{9} And said: O Lord, Lord, almighty king, for all things are in thy power, and there is none that can resist thy will, if thou determine to save Israel.

\textsuperscript{10} Thou hast made heaven and earth and all things that are under the cope of heaven.

\textsuperscript{11} Thou art Lord of all, and there is none that can resist thy majesty.
\end{englishwitness}

\subsection{Collect}
\label{proper-collect}
\properlocus{\latin{Omnipotens sempiterne Deus, qui abundantia}.}
Roman Missal, Twenty-seventh Sunday in Ordinary Time, Collect. The approved English and full Latin prayer are not reproduced. The prayer's petition is discussed above; this locator supplies no replacement text.

\subsection{First reading}
\label{proper-first-reading}
\properlocus{Isaiah 5:1--7.}

\begin{englishwitness}{Douay--Rheims Challoner}
\textsuperscript{1} I will sing to my beloved the canticle of my cousin concerning his vineyard. My beloved had a vineyard on a hill in a fruitful place.

\textsuperscript{2} And he fenced it in, and picked the stones out of it, and planted it with the choicest vines, and built a tower in the midst thereof, and set up a winepress therein: and he looked that it should bring forth grapes, and it brought forth wild grapes.

\textsuperscript{3} And now, O ye inhabitants of Jerusalem, and ye men of Juda, judge between me and my vineyard.

\textsuperscript{4} What is there that I ought to do more to my vineyard, that I have not done to it? was it that I looked that it should bring forth grapes, and it hath brought forth wild grapes?

\textsuperscript{5} And now I will shew you what I will do to my vineyard. I will take away the hedge thereof, and it shall be wasted: I will break down the wall thereof, and it shall be trodden down.

\textsuperscript{6} And I will make it desolate: it shall not be pruned, and it shall not be digged: but briers and thorns shall come up: and I will command the clouds to rain no rain upon it.

\textsuperscript{7} For the vineyard of the Lord of hosts is the house of Israel: and the man of Juda, his pleasant plant: and I looked that he should do judgment, and behold iniquity: and do justice, and behold a cry.
\end{englishwitness}

\subsection{Responsorial Psalm}
\label{proper-responsorial-psalm}
\properlocus{Psalm 80:9,12,13--14,15--16,19--20; Douay--Rheims / Vulgate Psalm 79.}
The response is Isaiah 5:7a, whose full verse appears in the first reading above. Its appointed clause identifies the vineyard as the house of Israel; the full verse is not repeated as a response. The following are the selected psalm verses in study English.
\begin{englishwitness}{Douay--Rheims Challoner}
\textsuperscript{9} Thou hast brought a vineyard out of Egypt: thou hast cast out the Gentiles and planted it.

\textsuperscript{12} It stretched forth its branches unto the sea, and its boughs unto the river.

\textsuperscript{13} Why hast thou broken down the hedge thereof, so that all they who pass by the way do pluck it?

\textsuperscript{14} The boar out of the wood hath laid it waste: and a singular wild beast hath devoured it.

\textsuperscript{15} Turn again, O God of hosts, look down from heaven, and see, and visit this vineyard:

\textsuperscript{16} And perfect the same which thy right hand hath planted: and upon the son of man whom thou hast confirmed for thyself.

\textsuperscript{19} And we depart not from thee, thou shalt quicken us: and we will call upon thy name.

\textsuperscript{20} O Lord God of hosts, convert us and shew thy face, and we shall be saved.
\end{englishwitness}

\subsection{Second reading}
\label{proper-second-reading}
\properlocus{Philippians 4:6--9.}

\begin{englishwitness}{Douay--Rheims Challoner}
\textsuperscript{6} Be nothing solicitous: but in every thing, by prayer and supplication, with thanksgiving, let your petitions be made known to God.

\textsuperscript{7} And the peace of God, which surpasseth all understanding, keep your hearts and minds in Christ Jesus.

\textsuperscript{8} For the rest, brethren, whatsoever things are true, whatsoever modest, whatsoever just, whatsoever holy, whatsoever lovely, whatsoever of good fame, if there be any virtue, if any praise of discipline: think on these things.

\textsuperscript{9} The things which you have both learned and received and heard and seen in me, these do ye: and the God of peace shall be with you.
\end{englishwitness}

\subsection{Gospel acclamation}
\label{proper-acclamation}
\properlocus{Alleluia; cf. John 15:16.}
This is the full source verse for study. The Lectionary acclamation adapts it and does not appoint every clause of the verse.
\begin{englishwitness}{Douay--Rheims Challoner}
\textsuperscript{16} You have not chosen me: but I have chosen you; and have appointed you, that you should go and should bring forth fruit; and your fruit should remain: that whatsoever you shall ask of the Father in my name, he may give it you.
\end{englishwitness}

\subsection{Gospel}
\label{proper-gospel}
\properlocus{Matthew 21:33--43.}

\begin{englishwitness}{Douay--Rheims Challoner}
\textsuperscript{33} Hear ye another parable. There was a man, an householder, who planted a vineyard and made a hedge round about it and dug in it a press and built a tower and let it out to husbandmen and went into a strange country.

\textsuperscript{34} And when the time of the fruits drew nigh, he sent his servants to the husbandmen that they might receive the fruits thereof.

\textsuperscript{35} And the husbandmen laying hands on his servants, beat one and killed another and stoned another.

\textsuperscript{36} Again he sent other servants, more than the former; and they did to them in like manner.

\textsuperscript{37} And last of all he sent to them his son, saying: They will reverence my son.

\textsuperscript{38} But the husbandmen seeing the son, said among themselves: This is the heir: come, let us kill him, and we shall have his inheritance.

\textsuperscript{39} And taking him, they cast him forth out of the vineyard and killed him.

\textsuperscript{40} When therefore the lord of the vineyard shall come, what will he do to those husbandmen?

\textsuperscript{41} They say to him: He will bring those evil men to an evil end and will let out his vineyard to other husbandmen that shall render him the fruit in due season.

\textsuperscript{42} Jesus saith to them: Have you never read in the Scriptures: The stone which the builders rejected, the same is become the head of the corner? By the Lord this has been done; and it is wonderful in our eyes.

\textsuperscript{43} Therefore I say to you that the kingdom of God shall be taken from you and shall be given to a nation yielding the fruits thereof.
\end{englishwitness}

\subsection{Prayer over the Offerings}
\label{proper-offerings}
\properlocus{\latin{Suscipe, quaesumus, Domine}.}
Roman Missal, Twenty-seventh Sunday in Ordinary Time, Prayer over the Offerings. The exact approved-English exemplar remains uncollated. The analysis uses the identified 2002 Latin digital reproduction; neither that prayer body nor a reconstructed English version is printed.

\subsection{Communion: Lamentations alternative}
\label{proper-communion-lamentations}
\properlocus{Lamentations 3:25.}
First appointed Missal alternative. The retained Bible witness includes its alphabetic stanza marker, Teth; this is study Scripture, not an instruction to recite that marker in the antiphon.
\begin{englishwitness}{Douay--Rheims Challoner}
\textsuperscript{25} Teth. The Lord is good to them that hope in him, to the soul that seeketh him.
\end{englishwitness}

\subsection{Communion: Corinthians alternative}
\label{proper-communion-corinthians}
\properlocus{Cf. 1 Corinthians 10:17; 10:16 supplies the chalice context.}
Second appointed Missal alternative. The biblical verses explain its basis; the Missal adaptation is not identical with this two-verse extract. This alternative and the preceding one are not combined into one enacted antiphon.
\begin{englishwitness}{Douay--Rheims Challoner}
\textsuperscript{16} The chalice of benediction which we bless, is it not the communion of the blood of Christ? And the bread which we break, is it not the partaking of the body of the Lord?

\textsuperscript{17} For we, being many, are one bread, one body: all that partake of one bread.
\end{englishwitness}

\subsection{Prayer after Communion}
\label{proper-after-communion}
\properlocus{\latin{Concede nobis, omnipotens Deus}.}
Roman Missal, Twenty-seventh Sunday in Ordinary Time, Prayer after Communion. The approved English and full Latin are not reproduced. Its petition for a life transformed through reception is explained in the commentary.
